\documentclass[11pt]{article}
\usepackage[margin=0.9in]{geometry}
\usepackage{amsmath,booktabs,tabularx,graphicx,enumitem,xcolor,hyperref}
\graphicspath{{../outputs/causal_v2/runs/figures/}{../outputs/causal_v2/results/}}
\hypersetup{colorlinks=true,linkcolor=blue,urlcolor=blue,citecolor=blue}
% Generated from the audited causal_v2 result tables.
\newcommand{\TestWindows}{34,161}
\newcommand{\CachedWindows}{157,521}
\newcommand{\CleanMinADE}{0.330}
\newcommand{\CleanMinFDE}{0.578}
\newcommand{\CleanMeanADE}{0.855}
\newcommand{\CleanEnergy}{0.540}
\newcommand{\NoiseMeanDelta}{+2.178}
\newcommand{\MissingMeanDelta}{+0.092}
\newcommand{\NoiseMeanInterval}{[1.757, 2.586]}
\newcommand{\MissingMeanInterval}{[0.032, 0.138]}
\newcommand{\OneSampleADE}{0.857}
\newcommand{\CleanSpeed}{9.4\%}
\newcommand{\CleanAccel}{13.8\%}
\newcommand{\CleanJerk}{13.4\%}
\newcommand{\CleanClose}{87.2\%}
\newcommand{\RuntimePre}{30.8}
\newcommand{\RuntimePred}{7.4}
\newcommand{\RuntimeTotal}{38.7}
\newcommand{\FailureADE}{4.767}
\newcommand{\FailureFold}{ETH}
\newcommand{\FoldRows}{
ETH & 364 & 0.600 & 0.021 & 1.037 & 1.366 & 0.935 \\
Hotel & 1,197 & 0.276 & 0.023 & 0.467 & 0.857 & 0.467 \\
Univ & 24,334 & 0.312 & 0.006 & 0.586 & 0.722 & 0.523 \\
Zara1 & 2,356 & 0.238 & 0.000 & 0.397 & 0.709 & 0.399 \\
Zara2 & 5,910 & 0.222 & 0.016 & 0.401 & 0.621 & 0.378 \\
}
\newcommand{\SlideFoldRows}{
ETH & 0.600 & 1.366 & 0.935 \\
Hotel & 0.276 & 0.857 & 0.467 \\
Univ & 0.312 & 0.722 & 0.523 \\
Zara1 & 0.238 & 0.709 & 0.399 \\
Zara2 & 0.222 & 0.621 & 0.378 \\
}
\newcommand{\ConditionRows}{
Clean & 0.330 & 0.578 & 0.855 & 1.823 & 0.540 & 0.991 \\
Noise 0.05 m & 0.485 & 0.744 & 1.482 & 2.978 & 0.768 & 1.966 \\
Noise 0.10 m & 0.683 & 0.984 & 2.094 & 4.111 & 1.029 & 2.851 \\
Noise 0.20 m & 1.054 & 1.426 & 3.033 & 5.804 & 1.510 & 4.051 \\
Missing one & 0.373 & 0.628 & 0.895 & 1.848 & 0.592 & 0.950 \\
Missing two & 0.437 & 0.697 & 0.947 & 1.881 & 0.662 & 0.894 \\
}

\title{Trajectron++ on ETH/UCY\\\large Independent Pedestrian Forecasting with Causal Inputs}
\author{LI Junhao \quad COMP5331 Group 7}
\date{Fall 2026}
\begin{document}
\maketitle

\begin{abstract}
This report presents the independent Trajectron++ contribution to the Group 7
comparison of six pedestrian forecasting models. The trajectory-only model
uses graph context, a discrete conditional variational autoencoder, and a
velocity integrator. We removed look-ahead through derived motion features
and verified that changes after the prediction anchor leave inputs and
forecasts invariant. The corrected matrix contains 15 independently trained
fold--seed runs, 90 test bundles, and \TestWindows{} unique test windows.
Clean scene-macro minADE, minFDE, and Energy Score are \CleanMinADE{} m,
\CleanMinFDE{} m, and \CleanEnergy{} m, respectively. The evaluation reports
the full sampling and robustness profile with traceable inputs, selected
checkpoints, and a model-independent evaluator.
\end{abstract}

\section{Task and scope}
Given eight metric positions per target at 0.4 s intervals and eligible
neighbor histories, forecast twelve positions over the following 4.8 s.
The output is twenty equally weighted trajectories per target. This is the
pedestrian, trajectory-only interpretation of Trajectron++\cite{tpp}.
It uses ETH, Hotel, Univ, Zara1 and Zara2\cite{eth,ucy}. The independent
implementation and shared interfaces address Raymond's requirements for
independent model development, common data and metrics, documented deviations,
and reproducible comparison. Results here describe one model contribution.
The group ranking requires matching evidence from the five other owners.

\section{Independent model}
At each historical step the target uses relative position, velocity divided by
two, and acceleration. Each neighbor must have all eight history points.
The graph masks each pedestrian edge outside 3 m of the target and sums states
within the single edge type. Target and edge LSTMs encode the histories.
An attention module acts over the available edge type, and the resulting
context parameterizes a categorical prior over 25 latent classes.

During training, a bidirectional future-position encoder supplies the CVAE
posterior. The loss combines normalized velocity negative log likelihood and
a KL term with annealing and a 0.07 floor. An autoregressive GRU predicts the
mean, scales and correlation of a bivariate Gaussian velocity at each step.
At inference the model samples the prior and decoder with history-only
context. The single integrator computes
\[
 \hat p_{t+h}=p_t^{\mathrm{filled}} + 0.4\sum_{j=1}^{h} \hat v_{t+j},
\]
where velocity is restored to meters per second before integration. This
produces position sequences in the original metric frame. Neighbor interactions
enter through historical context. Individual forecast samples carry no joint
collision constraint.

\section{Causal input and experimental design}
\paragraph{Leakage correction.}
The earlier pipeline computed centered derivatives along complete tracks and
then sliced windows. The velocity at the anchor could include the first future
point and the acceleration could include the second. The corrected adapter
derives target and neighbor motion inside each observed window with backward
differences. The first velocity and first two accelerations are zero.
The training velocity labels use the true anchor and future position
differences and have a separate path from inference. Cache entries bind the
causal protocol to the input file and ordered manifest hashes.

\paragraph{Verification.}
For every condition the regression changes future target or neighbor
positions while holding history fixed. The feature tensors and same-seed
predictions must remain exactly equal. The data audit also checks all
\CachedWindows{} cached windows against immediate construction, rejects a
legacy cache, and checks all 15 split manifests against the locked original
ordering. The former outputs are archived separately from this matrix.

\paragraph{Splits and selection.}
Hold out one of five scenes for testing and use temporal development splits
within the remaining sources. The boundary excludes 20 steps on either side
of the 80\% split point. All runs use seeds 0, 1 and 2, Adam at $10^{-3}$,
gradient clipping at 1, 100 epochs and training batches of 256.
Evaluate the epoch 10 through 100 checkpoints at intervals of ten on clean
validation data with $K=20$. Select the smallest Energy Score, then smaller
meanADE, then earlier epoch. Freeze the selection before any test prediction.

\paragraph{Perturbations.}
Conditions are clean, Gaussian position noise at 0.05, 0.10 and 0.20 m, and
missing the last one or two target observations. Missing values are filled
with constant velocity estimated from the two latest visible points.
Noise is applied to the target history only and sampled by NumPy
\texttt{RandomState} using $100003s + 1009i$ for training seed $s$ and zero-based
manifest row $i$. Coordinates are traversed in time order, x before y.
Neighbor positions remain available up to the anchor. Future labels are
identical across conditions. Exported masks distinguish measured and filled
positions.

\section{Evaluation and results}
For sample $k$, ADE averages Euclidean position error over twelve steps and
FDE measures endpoint error. Min metrics independently minimize each error
over $K=20$ samples. Mean metrics average over the twenty samples.
Energy Score uses the root-mean-square Euclidean distance over a trajectory
and subtracts half the average distance over distinct sample pairs from the
average sample-to-truth distance\cite{energy}. Diversity reports that mean
pair distance. Lower displacement error or Energy Score is preferred.
Diversity describes spread and has no universal optimum.

Compute metrics per window, average within each held-out fold and seed, then
average the three seeds within each fold and give each fold equal weight.
This prevents the large Univ fold from dominating the headline result.
Seed standard deviations are reported per fold. Paired perturbation changes
are computed within folds and the bootstrap resamples the five folds, rather
than treating overlapping windows as independent. Intervals are descriptive.

\begin{table}[htbp]
\centering
\caption{Clean holdout results in meters. The minADE standard deviation is across three training seeds.}
\begin{tabular}{lrrrrrr}
\toprule
Held-out & Windows & minADE & SD & minFDE & meanADE & ES \\
\midrule
\FoldRows
\bottomrule
\end{tabular}
\end{table}

\begin{table}[htbp]
\centering
\caption{Equal-weight fold macro results in meters. Lower error and ES are preferred. Div. denotes sample diversity.}
\begin{tabular}{lrrrrrr}
\toprule
Condition & minADE & minFDE & meanADE & meanFDE & ES & Div. \\
\midrule
\ConditionRows
\bottomrule
\end{tabular}
\end{table}

For clean inputs the gap between minADE (\CleanMinADE{} m) and meanADE
(\CleanMeanADE{} m) indicates that sample coverage and typical sampled error
describe different aspects of the forecast. The paired meanADE differences
for 0.20 m noise and two missing observations are \NoiseMeanDelta{} m and
\MissingMeanDelta{} m, with 95\% scene bootstrap intervals of
\NoiseMeanInterval{} m and \MissingMeanInterval{} m, respectively.
These changes describe the registered target-history corruption and fill
protocol. The same clean sample stream yields \OneSampleADE{} m with its
first trajectory and \CleanMinADE{} m best-of-twenty. The sample-budget table
also reports the intermediate budgets of five and ten.

\section{Motion, proximity and failure cases}
The evaluator derives predicted speed, acceleration and jerk from the last
input position and predicted positions. For each fold, the 99th percentiles
of training-track motion define diagnostic thresholds. A window is flagged
when the 99th percentile across its sampled motion exceeds the threshold.
The clean macro flag fractions are \CleanSpeed{} for speed, \CleanAccel{}
for acceleration, and \CleanJerk{} for jerk. Heading diagnostics exclude
steps below 0.05 m/s. Turning strata use ground-truth future heading and
density strata use anchor-time neighbors, solely for post-hoc analysis.

For each frame with at least two eligible targets, the evaluator matches
equal sample indices, takes the minimum separation over target pairs and
forecast times, and counts how often it falls below the training fifth
percentile of contemporaneous separation. The clean macro close-approach
fraction is \CleanClose{}. Equal-index matching is a descriptive construction
over separately sampled forecasts. Its fraction is neither a calibrated
joint collision probability nor a physical safety guarantee.

\begin{figure}[htbp]
 \centering\includegraphics[width=\linewidth]{eth_presentation.png}
 \caption{Fixed ETH example, seed 0. Black shows observations, dashed black
 shows causal fill, orange shows the ground truth, and blue shows twenty
 predicted trajectories and endpoints. Panels share their coordinate scale.}
\end{figure}

The archive provides fixed cases from each source and failure selections for
large clean error, deterioration under 0.20 m noise, deterioration with two
missing steps, and simultaneous motion flags. The largest clean meanADE case
in the seed-0 failure list has meanADE \FailureADE{} m on \FailureFold{}.
This rule exposes difficult windows alongside the fixed examples. The
qualitative plots complement the scene-level summaries and do not isolate
effects attributable to specific architecture components.

\section{Runtime and reproducibility}
Use the first 128 target windows of each held-out manifest with batch 32 and
twenty samples. For each selected model, warm up three times and measure ten
repetitions with CUDA synchronization at stage boundaries. Preprocessing
starts with in-memory neutral positions and includes window and graph
construction, collation and transfer to the GPU. Inference includes encoding,
sampling, decoding and metric integration. Total time also includes CPU
readback. Disk I/O and one-time model or data loading are outside the timed
boundary. The median across 15 per-run medians is \RuntimePre{} ms
preprocessing, \RuntimePred{} ms inference with sampling, and \RuntimeTotal{}
ms end to end. Hardware and per-repetition data are included with the archive.

The prediction index resolves each output to a fold, seed, condition, ordered
manifest hash, configuration hash, selected checkpoint hash and sample count.
The README includes locked data acquisition, environment setup, metric
fixtures, prediction, evaluation, matrix and report commands. The independent
path uses no upstream model, trainer or evaluator at runtime. The reference
code was consulted for algorithm details and pedestrian hyperparameters.

\section{Deviations and limits}
The paper supports multiple agent types, semantic maps and optional
ego-agent plan conditioning. This implementation uses one pedestrian edge
type, complete neighbor histories, a hard radius mask and trajectory inputs.
Window-local normalization and differentiation, the velocity NLL training
objective, and sampled integration are implementation choices. Group 7
specifies its own split, perturbations and $K=20$ budget. Published numbers
therefore provide context under their original protocol. Our benchmark covers
five scenes, so the results characterize these holdouts and the declared
perturbations.

\begin{thebibliography}{9}
\small
\setlength{\itemsep}{0pt}
\bibitem{tpp} T. Salzmann, B. Ivanovic, P. Chakravarty, and M. Pavone.
Trajectron++: Dynamically-Feasible Trajectory Forecasting With Heterogeneous
Data. ECCV, 2020. \href{https://arxiv.org/abs/2001.03093}{Paper}.
\bibitem{eth} S. Pellegrini, A. Ess, K. Schindler, and L. Van Gool.
You'll Never Walk Alone: Modeling Social Behavior for Multi-target Tracking.
ICCV, 2009.
\bibitem{ucy} A. Lerner, Y. Chrysanthou, and D. Lischinski. Crowds by
Example. Computer Graphics Forum, 26(3), 2007.
\bibitem{energy} T. Gneiting and A. E. Raftery. Strictly Proper Scoring Rules,
Prediction, and Estimation. JASA, 102(477), 2007.
\href{https://doi.org/10.1198/016214506000001437}{DOI}.
\end{thebibliography}
\end{document}
